\section*{Laplace Transformation Aufgaben mit L�sungen}\label{MA2-laplacetransformation-loesungen}

\ifintegriereaufgaben Erkl�rungen und Beispiele, die zum Verst�ndnis der L�sungswege hilfreich sind, finden Sie im Skript\wennreferenzvorhanden{MA2-laplacetransformation}{ (Abschnitt \abschnittsverweis{MA2-laplacetransformation})}. \fi

\begin{exercisebox}[Uneigentliche Integrale] %A.Z. Integralrechnung  % Papula I.V, Abschnitt 9, Aufgabe 1 und 2
	Bestimmen Sie den Wert der folgenden uneigentlichen Integrale, falls er existiert.
	\begin{eqnarray*}
		(a)\;\int_0^\infty e^{-x}\,\mathrm{d}x\quad
		(b)\;\int_0^\infty x\cdot e^{-x}\,\mathrm{d}x\quad
		(c)\;\int_{-\infty}^2 e^{x}\,\mathrm{d}x\quad
		(d)\;\int_{0}^\infty x^2\cdot e^{-2x}\,\mathrm{d}x
	\end{eqnarray*}
\vspace{0.3cm}
\textbf{L�sung:} \newline
Diese Aufgaben entsprechen den Aufgaben \(1\) und \(2\) zu Abschnitt V.\(9\) in Papula Band 1. Die Aufgabenstellungen finden Sie auf Seite \(564\) und detaillierte L�sungswege auf Seite \(798\).
\begin{itemize}
	\item[(a)] Papula 1.V.9, Aufg. 1.a) $\quad 1$
	\item[(b)] Papula 1.V.9, Aufg. 1.b) $\quad 1$
	\item[(c)] Papula 1.V.9, Aufg. 1.c) $\quad e^2$
	\item[(d)] Papula 1.V.9, Aufg. 2) $\quad 1/4$
\end{itemize}
\end{exercisebox}

\begin{exercisebox}[Uneigentliche Integrale] %A.Z. Integralrechnung  % Papula I.V, Abschnitt 9, Aufgabe 4
	Bestimmen Sie den Wert der folgenden uneigentlichen Integrale, falls er existiert.
	\begin{eqnarray*}
		(a)\;\int_{-1}^0 \frac{1}{\sqrt{1+x}}\,\mathrm{d}x\quad
		(b)\;\int_{-1}^{1} \frac{1}{x^2}\,\mathrm{d}x\quad
		(c)\;\int_{0}^{10} \frac{e^{x}}{e^x-1}\,\mathrm{d}x
	\end{eqnarray*}
\vspace{0.3cm}
\textbf{L�sung:} \newline
Diese Aufgabe entspricht Aufgabe \(4\) zu Abschnitt V.\(9\) in Papula Band 1. Die Aufgabenstellungen finden Sie auf Seite \(565\) und detaillierte L�sungswege auf Seite \(799\).
\begin{itemize}
	\item[(a)] $2$
	\item[(b)] divergent
	\item[(c)] divergent
\end{itemize}
\end{exercisebox}


\begin{exercisebox}[Laplace-Transformation] % Papula II VI, Abschnitt 1, Aufgabe 1 und 3
	Aufgabe war es zu den gegebenen Funktionen die Laplace-Transformierte durch explizite Berechnung zu bestimmen. Die Aufgaben finden Sie im Lehrbuch von Papula Band II, Abschnitt VI.\(1\) auf Seite \(676\) und detaillierte L�sungswege auf Seite \(766\):
	\begin{itemize}
	\item[(a)] Papula, Aufg. 1.b)  $\quad f(t) = 2t\cdot e^{-4t}$
	\item[(b)] Papula, Aufg. 1.c)  $\quad f(t) = e^{-\delta t} \cdot \sin(\omega t), \; \delta, \omega >0$
	\item[(c)] Papula, Aufg. 1.d)  $\quad f(t) = \mathrm{sinh}(at),\, a>0$
	\item[(d)] Papula, Aufg. 1.f)  $\quad f(t) = \sin^2(t)$
\item[(e)] Papula, Aufg. 3.a) \(\quad\) (Rechteckimpuls) F�r $a>0$
			\[ f(t) = \begin{cases} A, \quad & 0<t<a\,,\\
					-A, \quad & a< t < 2a\,,\\
			0, \quad \mathrm{sonst}\,. \end{cases}\]
		\item[(f)] Papula, Aufg. 3.b) \(\quad\) (Sinusimpuls) F�r $a>0$
			\[ f(t) = \begin{cases} A\cdot \sin\left( \frac{\pi\cdot t}{a}\right), \quad & 0\leq t \leq a\,,\\
			0, \quad \mathrm{sonst}\,. \end{cases}\]
	\end{itemize}
\end{exercisebox}


\begin{exercisebox}[Laplace-Transformation] % Papula II, VI, Abschnitt 2, Aufgabe 4
	Bestimmen Sie
	%mit Hilfe von
	%\[ \sin(t) \laplace \frac{1}{s^2+1} \]
	die Laplace-Transformierte von
	\[ f(t) = \sin(\omega t + \varphi), \; \omega,\varphi >0\,. \]
	%indem Sie zuerst eine Verschiebung und dann eine Skalierung verwenden.

\vspace{0.3cm}
\textbf{L�sung:} \newline
		\begin{eqnarray*}
	I &=& \int\limits_0^\infty \sin(\omega t +\varphi) e^{-st}\,\mathrm{d}t = \frac{1}{s}\sin(\varphi) + \frac{\omega}{s} \int\limits_{0}^\infty \cos(\omega t+\varphi) e^{-st}\,\mathrm{d}t \\ &=& \frac{1}{s}\sin(\varphi) + \frac{\omega}{s^2}\cos(\varphi) - \frac{\omega^2}{s^2} I  \\
	&\Rightarrow &  I \left( 1+\frac{\omega^2}{s^2}\right) = \frac1s\sin(\varphi) - \frac{\omega}{s^2}\cos(\varphi)
\end{eqnarray*}
Also erhalten wir
\[ \Rightarrow I = \mathcal L(f) = \frac{1}{s^2} \left( s \sin(\varphi) - \omega \cos(\varphi) \right)  \frac{s^2}{s^2 + \omega^2} = \frac{  s \sin(\varphi) - \omega \cos(\varphi)}{s^2 + \omega^2}\,.\]

Einen alternativen L�sungsweg, der von den S�tzen �ber die Laplace-Transformierte bei " Manipulationen vom Argument " (Beschleunigung und Verschiebung des Signals) Gebrauch macht, finden Sie im Lehrbuch von Papula Band II, Abschnitt VI.\(2\) finden Sie auf Seite \(767f\). Die Aufgabe ist Aufgabe \(4\) auf Seite \(678\).
\end{exercisebox}
